\chapter{线性代数：行列式和矩阵}

\section{行列式、伴随矩阵与矩阵运算}

\begin{problemblock}
\textbf{例 1.} 多项式
\[
f(x)=\begin{vmatrix}
x&x&1&2x\\
1&x&2&-1\\
2&1&x&1\\
2&-1&1&x
\end{vmatrix}
\]
中 \(x^3\) 项的系数为多少？
\end{problemblock}
\begin{solutionblock}
\analysis{这是行列式按参数展开求指定系数的问题。直接把行列式看成关于 \(x\) 的多项式，按行列式定义或用初等变换展开即可。计算得
\[
f(x)=x^4-5x^3-2x^2+13x+5.
\]
因此 \(x^3\) 项的系数为 \(-5\)。考研中若只求某一项系数，也可按排列乘积中取三个含 \(x\) 因子的项分类，但本题四阶行列式直接化简更稳。}
\answer{\(-5\)}
\examnote{含参行列式求某项系数时，既可整体算出多项式，也可按“选取含参元素个数”分类；四阶以内整体计算通常更省错。}
\end{solutionblock}

\begin{problemblock}
\textbf{例 2.} 设
\[
A=\begin{pmatrix}
a+1&b&3\\
a&\dfrac b2&1\\
1&1&2
\end{pmatrix},
\]
\(M_{ij}\) 为 \(a_{ij}\) 的余子式。若 \(|A|=-\dfrac12\)，且
\[
-M_{21}+M_{22}-M_{23}=0,
\]
则下列结论正确的是（\quad）：
\[
\begin{array}{ll}
\text{A. }a=0\text{ 或 }a=-\dfrac32, & \text{B. }a=0\text{ 或 }a=\dfrac32,\\[4pt]
\text{C. }b=1\text{ 或 }b=-\dfrac12, & \text{D. }b=-1\text{ 或 }b=\dfrac12.
\end{array}
\]
\end{problemblock}
\begin{solutionblock}
\analysis{先计算两个条件。由行列式展开可得
\[
|A|=-\frac{(2a-1)(b-2)}2.
\]
再算余子式：
\[
-M_{21}+M_{22}-M_{23}=a-b+1.
\]
所以条件化为
\[
-(2a-1)(b-2)/2=-\frac12,\qquad a-b+1=0.
\]
第二式给出 \(b=a+1\)。代入第一式：
\[
(2a-1)(a-1)=1
\quad\Longleftrightarrow\quad
2a^2-3a=0
\quad\Longleftrightarrow\quad
a=0\text{ 或 }a=\frac32.
\]
对应选项为 B。}
\answer{B；\(a=0\) 或 \(a=\dfrac32\)}
\examnote{余子式 \(M_{ij}\) 与代数余子式 \(A_{ij}=(-1)^{i+j}M_{ij}\) 不要混淆，本题写的是 \(M_{ij}\)。}
\end{solutionblock}

\begin{problemblock}
\textbf{例 3.} 设
\[
\begin{vmatrix}
a&b&c&d\\
1&1&2&2\\
x&y&z&w\\
2&2&1&1
\end{vmatrix}=6,
\]
求 \(A_{21}+A_{22}\)，其中 \(A_{ij}\) 为该行列式中元素的代数余子式。
\end{problemblock}
\begin{solutionblock}
\analysis{代数余子式 \(A_{2j}\) 与第二行元素无关。把第二行改成 \((1,1,0,0)\)，则新行列式等于
\[
A_{21}+A_{22}.
\]
记原行列式为 \(D=6\)。原第二行 \(R_2=(1,1,2,2)\)，第四行 \(R_4=(2,2,1,1)\)。由于
\[
(1,1,0,0)=-\frac13R_2+\frac23R_4,
\]
把第二行替换为上式后，含 \(R_4\) 的那一项行列式有两行相同而为 \(0\)，故
\[
A_{21}+A_{22}=-\frac13D=-2.
\]}
\answer{\(-2\)}
\examnote{求一行若干代数余子式的线性组合，常把对应行替换为组合系数行，这是非常高频的行列式技巧。}
\end{solutionblock}

\begin{problemblock}
\textbf{例 4.} 设 \(A\) 为三阶方阵，\(A_{ij}\) 是 \(A\) 中第 \(i\) 行第 \(j\) 列元素的代数余子式，
\[
B=\begin{pmatrix}0&0&0\\0&2&0\\0&0&0\end{pmatrix}.
\]
若 \(|A|=2,\ |A+B|=12\)，求 \(A_{22}\)。
\end{problemblock}
\begin{solutionblock}
\analysis{矩阵 \(B\) 只把 \(A\) 的 \(a_{22}\) 增加了 \(2\)。行列式关于某一个元素是一次函数，并且该元素的系数就是对应代数余子式 \(A_{22}\)。因此
\[
|A+B|=|A|+2A_{22}.
\]
代入 \(|A|=2,\ |A+B|=12\)，得
\[
12=2+2A_{22},
\]
所以 \(A_{22}=5\)。}
\answer{\(5\)}
\examnote{“只改一个元素”时，行列式增量等于“改变量 \(\times\) 对应代数余子式”。}
\end{solutionblock}

\begin{problemblock}
\textbf{例 5.} 设三阶矩阵 \(A=(a_{ij})\) 满足各行元素之和均为 \(2\)，且 \(|A|=\dfrac12\)。若矩阵 \(B(t)=(a_{ij}+t)\)，求 \(|B(1)|\)。
\end{problemblock}
\begin{solutionblock}
\analysis{记 \(\mathbf 1=(1,1,1)^T\)，则“各行和为 \(2\)”等价于
\[
A\mathbf 1=2\mathbf 1.
\]
又 \(B(1)=A+J\)，其中 \(J=\mathbf1\mathbf1^T\)。由矩阵行列式引理
\[
|A+\mathbf1\mathbf1^T|
=|A|\left(1+\mathbf1^TA^{-1}\mathbf1\right).
\]
由 \(A\mathbf1=2\mathbf1\) 得 \(A^{-1}\mathbf1=\dfrac12\mathbf1\)，故
\[
\mathbf1^TA^{-1}\mathbf1
=\mathbf1^T\frac12\mathbf1
=\frac32.
\]
于是
\[
|B(1)|=\frac12\left(1+\frac32\right)=\frac54.
\]}
\method{方法二：特征值观点}
\(A\mathbf1=2\mathbf1\)，而 \(J\) 在 \(\mathbf1\) 方向上的特征值为 \(3\)，在与 \(\mathbf1\) 正交的方向上为 \(0\)。对 \(A\) 乘上 \(I+A^{-1}J\) 后，只在 \(\mathbf1\) 方向贡献因子 \(1+\frac32\)，所以同样得到 \(|A+J|=\frac54\)。
\answer{\(\dfrac54\)}
\examnote{遇到“每行和相同”优先想到特征向量 \(\mathbf1\)，遇到“所有元素同时加 \(t\)”优先想到秩一矩阵 \(t\mathbf1\mathbf1^T\)。}
\end{solutionblock}

\begin{problemblock}
\textbf{例 6.} 计算行列式
\[
\begin{vmatrix}
1&0&2&-1\\
0&2&1&0\\
1&-1&0&1\\
1&2&3&4
\end{vmatrix}.
\]
\end{problemblock}
\begin{solutionblock}
\analysis{直接作初等行变换。用第一行消去第一列下面元素：
\[
R_3\leftarrow R_3-R_1,\qquad R_4\leftarrow R_4-R_1,
\]
得到
\[
\begin{vmatrix}
1&0&2&-1\\
0&2&1&0\\
0&-1&-2&2\\
0&2&1&5
\end{vmatrix}.
\]
沿第一列展开，只需计算后三阶行列式：
\[
\begin{vmatrix}
2&1&0\\
-1&-2&2\\
2&1&5
\end{vmatrix}
=2(-10-2)-1(-5-4)=-24+9=-15.
\]}
\answer{\(-15\)}
\examnote{数值行列式以行列变换为主，尽量先造零，再沿零多的行或列展开。}
\end{solutionblock}

\begin{problemblock}
\textbf{例 7.} 计算
\[
\begin{vmatrix}
a&0&-1&1\\
0&a&1&-1\\
-1&1&a&0\\
1&-1&0&a
\end{vmatrix}.
\]
\end{problemblock}
\begin{solutionblock}
\analysis{把矩阵看成 \(aE+K\)，其中
\[
K=\begin{pmatrix}
0&0&-1&1\\
0&0&1&-1\\
-1&1&0&0\\
1&-1&0&0
\end{pmatrix}.
\]
矩阵 \(K\) 的非零作用只在由 \((1,-1,0,0)^T\)、\((0,0,1,-1)^T\) 张成的二维子空间上；在其正交补上特征值为 \(0,0\)。在该二维子空间中可算得 \(K\) 的特征值为 \(2,-2\)。因此原矩阵特征值为
\[
a,\ a,\ a+2,\ a-2.
\]
行列式等于特征值乘积：
\[
a^2(a+2)(a-2)=a^2(a-2)(a+2).
\]}
\method{方法二：分块计算}
按 \(2\times2\) 分块写成
\[
\begin{pmatrix}
aI_2&C\\ C&aI_2
\end{pmatrix},\qquad
C=\begin{pmatrix}-1&1\\1&-1\end{pmatrix}.
\]
同加同减的分块变换可把行列式化为
\[
|aI_2+C|\,|aI_2-C|
=a(a-2)\cdot a(a+2).
\]
\answer{\(a^2(a-2)(a+2)\)}
\examnote{具有对称分块结构的行列式，常用“同加同减”把它化为两个低阶行列式。}
\end{solutionblock}

\begin{problemblock}
\textbf{例 8.} 计算
\[
\begin{vmatrix}
\lambda&-1&0&0\\
0&\lambda&-1&0\\
0&0&\lambda&-1\\
4&3&2&\lambda+1
\end{vmatrix}.
\]
\end{problemblock}
\begin{solutionblock}
\analysis{这是伴随矩阵型行列式。沿第一行递推或直接按第一列展开均可。设该行列式为 \(D(\lambda)\)。展开计算：
\[
\begin{aligned}
D(\lambda)
&=\lambda
\begin{vmatrix}
\lambda&-1&0\\
0&\lambda&-1\\
3&2&\lambda+1
\end{vmatrix}
+
\begin{vmatrix}
0&-1&0\\
0&\lambda&-1\\
4&2&\lambda+1
\end{vmatrix}\\
&=\lambda\bigl(\lambda^3+\lambda^2+2\lambda+3\bigr)+4\\
&=\lambda^4+\lambda^3+2\lambda^2+3\lambda+4.
\end{aligned}
\]}
\answer{\(\lambda^4+\lambda^3+2\lambda^2+3\lambda+4\)}
\examnote{上 Hessenberg 形行列式常用递推；不要把最后一行的 \(4,3,2,\lambda+1\) 顺序看反。}
\end{solutionblock}

\begin{problemblock}
\textbf{例 9.} 计算
\[
\begin{vmatrix}
1&2&3&4\\
1&2^2&3^2&4^2\\
1&2^3&3^3&4^3\\
9&8&7&6
\end{vmatrix}.
\]
\end{problemblock}
\begin{solutionblock}
\analysis{把每列看成由 \(t=1,2,3,4\) 代入向量
\[
(1,\ t,\ t^2,\ 10-t)^T
\]
得到。第 4 行 \(10-t\) 是 \(1,t,t^2\) 的线性组合中的一次表达。更直接地，对列作 Vandermonde 型化简，计算可得
\[
\begin{vmatrix}
1&2&3&4\\
1&4&9&16\\
1&8&27&64\\
9&8&7&6
\end{vmatrix}
=-120.
\]
也可将第 4 行写成 \(10\cdot(1,1,1,1)-(1,2,3,4)\)，再用 Vandermonde 行列式化到标准形式。}
\answer{\(-120\)}
\examnote{列参数分别为 \(1,2,3,4\) 的行列式，优先往 Vandermonde 行列式靠。}
\end{solutionblock}

\begin{problemblock}
\textbf{例 10.} 计算行列式
\[
\begin{vmatrix}
1+2a&a+2a^2&a^2+2a^3&2+a^3\\
1+2b&b+2b^2&b^2+2b^3&2+b^3\\
1+2c&c+2c^2&c^2+2c^3&2+c^3\\
1+2d&d+2d^2&d^2+2d^3&2+d^3
\end{vmatrix}.
\]
\end{problemblock}
\begin{solutionblock}
\analysis{每一行由 \(t=a,b,c,d\) 代入
\[
(1+2t,\ t+2t^2,\ t^2+2t^3,\ 2+t^3)
\]
得到。它等于标准基 \((1,t,t^2,t^3)\) 右乘一个列变换矩阵。按列写
\[
\begin{aligned}
1+2t&=1+2t,\\
t+2t^2&=t+2t^2,\\
t^2+2t^3&=t^2+2t^3,\\
2+t^3&=2+t^3.
\end{aligned}
\]
对应列变换矩阵的行列式为 \(-15\)。因此原行列式等于
\[
-15
\begin{vmatrix}
1&a&a^2&a^3\\
1&b&b^2&b^3\\
1&c&c^2&c^3\\
1&d&d^2&d^3
\end{vmatrix}.
\]
标准 Vandermonde 行列式为
\[
\prod_{i<j}(t_j-t_i)
=(b-a)(c-a)(d-a)(c-b)(d-b)(d-c).
\]}
\answer{\(-15(b-a)(c-a)(d-a)(c-b)(d-b)(d-c)\)}
\examnote{“每行由同一组多项式代入不同字母”时，答案一定是常数倍 Vandermonde。关键是算准列变换常数。}
\end{solutionblock}

\begin{problemblock}
\textbf{例 11.} 设四阶矩阵
\[
A=(\alpha,\gamma_1,\gamma_2,\gamma_3),\qquad
B=(\beta,\gamma_2,\gamma_3,\gamma_1),
\]
且 \(|A|=a,\ |B|=b\)。求 \(|A+B|\)。
\end{problemblock}
\begin{solutionblock}
\analysis{按列相加：
\[
A+B=(\alpha+\beta,\ \gamma_1+\gamma_2,\ \gamma_2+\gamma_3,\ \gamma_3+\gamma_1).
\]
利用行列式关于列的多重线性展开。后三列从每个括号各取一项，若出现两个相同的 \(\gamma_i\)，该项为 \(0\)。非零项只可能取出 \((\gamma_1,\gamma_2,\gamma_3)\) 或其循环排列。再看第一列取 \(\alpha\) 或 \(\beta\)。整理得
\[
|A+B|=2|\alpha,\gamma_1,\gamma_2,\gamma_3|
+2|\beta,\gamma_2,\gamma_3,\gamma_1|.
\]
其中 \((\gamma_1,\gamma_2,\gamma_3)\mapsto(\gamma_2,\gamma_3,\gamma_1)\) 是三循环，符号为正，所以
\[
|A+B|=2a+2b=2(a+b).
\]}
\answer{\(2(a+b)\)}
\examnote{多列同时相加的行列式不要硬展开全部 \(2^4\) 项，先用“重复列为零”删掉绝大多数项。}
\end{solutionblock}

\begin{problemblock}
\textbf{例 12.} \(A\) 是三阶矩阵，\(\alpha\) 是三维列向量，矩阵
\[
P=(\alpha,A\alpha,A^2\alpha)
\]
可逆，且
\[
A^3\alpha=3A\alpha-2A^2\alpha.
\]
求 \(|A+E|\)。
\end{problemblock}
\begin{solutionblock}
\analysis{在基
\[
\alpha,\quad A\alpha,\quad A^2\alpha
\]
下考察线性变换 \(A\)。有
\[
A\alpha=A\alpha,\qquad A(A\alpha)=A^2\alpha,\qquad
A(A^2\alpha)=3A\alpha-2A^2\alpha.
\]
因此 \(A\) 在该基下的矩阵为
\[
C=\begin{pmatrix}
0&0&0\\
1&0&3\\
0&1&-2
\end{pmatrix}.
\]
相似矩阵的 \(|A+E|\) 相同，所以
\[
|A+E|=|C+E|
=\begin{vmatrix}
1&0&0\\
1&1&3\\
0&1&-1
\end{vmatrix}
=-4.
\]}
\answer{\(-4\)}
\examnote{给出 \((\alpha,A\alpha,A^2\alpha)\) 可逆时，通常就是要求在这组基下写出 \(A\) 的表示矩阵。}
\end{solutionblock}

\begin{problemblock}
\textbf{例 13.} 设 \(A\) 为三阶矩阵。若
\[
|E-A|=|E+A|=|2E-A|=0,
\]
则 \(|3E-A|=(\quad)\)：
\[
\text{A. }-2,\qquad \text{B. }-8,\qquad \text{C. }8,\qquad \text{D. }11.
\]
\end{problemblock}
\begin{solutionblock}
\analysis{从特征值角度看，\(|\lambda E-A|=0\) 表示 \(\lambda\) 是 \(A\) 的特征值。条件
\[
|E-A|=0,\qquad |E+A|=0,\qquad |2E-A|=0
\]
说明 \(1,-1,2\) 均为 \(A\) 的特征值。三阶矩阵共有三个特征值，因此
\[
|3E-A|=(3-1)(3-(-1))(3-2)=2\cdot4\cdot1=8.
\]}
\answer{C；\(8\)}
\examnote{三阶矩阵给出三个不同的 \(|\lambda E-A|=0\)，就已经确定全部特征值。}
\end{solutionblock}

\begin{problemblock}
\textbf{例 14.} 设 \(A\) 为三阶矩阵。若
\[
|E-A|=1,\qquad |-E-A|=-1,\qquad |2E-A|=2,
\]
则 \(|3E-A|=(\quad)\)：
\[
\text{A. }-2,\qquad \text{B. }-8,\qquad \text{C. }8,\qquad \text{D. }11.
\]
\end{problemblock}
\begin{solutionblock}
\analysis{令
\[
p(t)=|tE-A|.
\]
因为 \(A\) 为三阶矩阵，\(p(t)\) 是首项系数为 \(1\) 的三次多项式，设
\[
p(t)=t^3+ut^2+vt+w.
\]
由题设
\[
p(1)=1,\qquad p(-1)=-1,\qquad p(2)=2.
\]
代入求得
\[
u=-2,\qquad v=0,\qquad w=2,
\]
故
\[
p(t)=t^3-2t^2+2.
\]
于是
\[
|3E-A|=p(3)=27-18+2=11.
\]}
\answer{D；\(11\)}
\examnote{\(|tE-A|\) 的首项一定是 \(t^n\)。这类题本质是用若干点值恢复特征多项式。}
\end{solutionblock}

\begin{problemblock}
\textbf{例 15.} 设向量
\[
\alpha_1=(a,1,-1)^T,\qquad
\alpha_2=(1,1,a)^T,\qquad
\alpha_3=(1,a,-1)^T.
\]
若 \(\alpha_1,\alpha_2,\alpha_3\) 线性相关，且其中任意两个向量均线性无关，则（\quad）。原题选项为
\[
\begin{array}{ll}
\text{A. }a=1,\ b\ne-1, & \text{B. }a=1,\ b=-1,\\
\text{C. }a\ne-2,\ b=2, & \text{D. }a=-2,\ b=2.
\end{array}
\]
\end{problemblock}
\begin{solutionblock}
\analysis{本题原 PDF 的向量 \(\alpha_2\) 第三个分量印成 \(a\)，但选项含 \(b\)。若严格按印出的 \((1,1,a)^T\)，则三向量线性相关只在 \(a=1\) 时发生，此时 \(\alpha_1=\alpha_3\)，不满足“任意两个线性无关”，四个选项都无法成立。结合选项可知命题应为
\[
\alpha_2=(1,1,b)^T.
\]
按这一意图，令三向量为列构成矩阵
\[
M=\begin{pmatrix}
a&1&1\\
1&1&a\\
-1&b&-1
\end{pmatrix}.
\]
三向量线性相关等价于
\[
|M|=-(a-1)(ab+b+2)=0.
\]
若 \(a=1\)，则
\[
\alpha_1=(1,1,-1)^T,\qquad \alpha_3=(1,1,-1)^T,
\]
任意两个无关条件被破坏，所以不能取 \(a=1\)。于是
\[
ab+b+2=0\quad\Longleftrightarrow\quad b(a+1)=-2.
\]
逐项检查选项，只有 \(a=-2,\ b=2\) 满足，并且此时任意两列均不成比例。}
\answer{按原 PDF 印出的题面无正确选项；若将 \(\alpha_2\) 第三个分量按选项意图修正为 \(b\)，则选 D，且 \(a=-2,\ b=2\)。}
\examnote{遇到“题面参数与选项参数不一致”时，不要硬套。应先按印刷题面检查是否有解，再说明若按命题意图修正后对应的选项。}
\end{solutionblock}

\begin{problemblock}
\textbf{例 16.} 已知
\[
A=\begin{pmatrix}1&1\\0&1\end{pmatrix},\qquad
B=\begin{pmatrix}1&-1\\0&1\end{pmatrix},
\]
\(C,D\) 均为二阶矩阵，则下列结论正确的是（\quad）：
\[
\begin{array}{ll}
\text{A. }\left|\begin{matrix}A&C\\B&D\end{matrix}\right|=|AD-BC|,&
\text{B. }\left|\begin{matrix}A&C\\B&D\end{matrix}\right|=|A|\,|D-CA^{-1}B|,\\[6pt]
\text{C. }|E-AB|\ne |E-BA|,&
\text{D. }|D-CA^{-1}B|\,|A|\ne |DA-CB|.
\end{array}
\]
\end{problemblock}
\begin{solutionblock}
\analysis{一般分块行列式不能随意写成 \(|AD-BC|\)，但本题的 \(A,B\) 有特殊结构：
\[
AB=BA=E.
\]
且 \(A\) 可逆。由 Schur 公式
\[
\left|\begin{matrix}A&C\\B&D\end{matrix}\right|
=|A|\cdot|D-BA^{-1}C|.
\]
因 \(A^{-1}=B\)，且 \(AB=E\)，有
\[
|A|\,|D-BA^{-1}C|
=|A(D-BA^{-1}C)|
=|AD-ABA^{-1}C|
=|AD-BC|.
\]
因此
\[
\left|\begin{matrix}A&C\\B&D\end{matrix}\right|=|AD-BC|.
\]
选项 B 把 Schur 补的顺序写错，一般不成立；选项 C 中恒有 \(|E-AB|=|E-BA|\)；选项 D 与 A 的正确等式矛盾。}
\answer{A}
\examnote{分块行列式公式要满足可交换或可逆 Schur 补条件，矩阵乘法顺序不能随意换。}
\end{solutionblock}

\begin{problemblock}
\textbf{例 17.} 设 \(A,B\) 为 \(n\) 阶矩阵，记 \(r(X)\) 为矩阵 \(X\) 的秩，\((X,Y)\) 表示分块矩阵，则（\quad）：
\[
\begin{array}{ll}
\text{A. }r(A,AB)=r(A), & \text{B. }r(A,BA)=r(A),\\
\text{C. }r(A,B)=\max\{r(A),r(B)\}, & \text{D. }r(A,B)=r(A^T,B^T).
\end{array}
\]
\end{problemblock}
\begin{solutionblock}
\analysis{矩阵 \(AB\) 的每一列都是 \(A\) 的列向量的线性组合，因此把 \(AB\) 拼在 \(A\) 后面不会增加列空间：
\[
r(A,AB)=r(A).
\]
所以 A 正确。B 不一定成立，因为 \(BA\) 的列空间未必包含在 \(A\) 的列空间内。C 也不正确，两个矩阵的列空间可能相互独立，秩可大于二者秩的最大值。D 中 \((A^T,B^T)\) 与 \((A,B)\) 的形状和列空间关系不同，一般不等秩。}
\answer{A}
\examnote{\((A,AB)\) 的关键是“右乘不改变列向量来自 \(A\) 列空间”的事实。}
\end{solutionblock}

\begin{problemblock}
\textbf{例 18.} 设 \(A,B\) 为 \(n\) 阶实矩阵，下列结论不成立的是（\quad）：
\[
\begin{array}{ll}
\text{A. }r\begin{pmatrix}A&O\\O&A^TA\end{pmatrix}=2r(A),&
\text{B. }r\begin{pmatrix}A&AB\\O&A^T\end{pmatrix}=2r(A),\\[8pt]
\text{C. }r\begin{pmatrix}A&BA\\O&AA^T\end{pmatrix}=2r(A),&
\text{D. }r\begin{pmatrix}A&O\\BA&A^T\end{pmatrix}=2r(A).
\end{array}
\]
\end{problemblock}
\begin{solutionblock}
\analysis{逐项看。实矩阵满足
\[
r(A^TA)=r(AA^T)=r(A),
\]
所以 A 正确。B 中右乘可逆分块矩阵
\[
\begin{pmatrix}I&-B\\O&I\end{pmatrix}
\]
得
\[
\begin{pmatrix}A&AB\\O&A^T\end{pmatrix}
\begin{pmatrix}I&-B\\O&I\end{pmatrix}
=\begin{pmatrix}A&O\\O&A^T\end{pmatrix},
\]
故 B 正确。D 中左乘可逆分块矩阵
\[
\begin{pmatrix}I&O\\-B&I\end{pmatrix}
\]
可化为
\[
\begin{pmatrix}A&O\\O&A^T\end{pmatrix},
\]
故 D 正确。C 的右上角为 \(BA\)，一般不能用同样的列变换消去。反例可取
\[
A=\begin{pmatrix}0&-1\\0&1\end{pmatrix},\qquad
B=\begin{pmatrix}-1&0\\1&1\end{pmatrix},
\]
则 \(r(A)=1\)，但
\[
r\begin{pmatrix}A&BA\\O&AA^T\end{pmatrix}=3\ne2.
\]}
\answer{C}
\examnote{这类秩题的核心是能否用可逆分块初等变换消掉非对角块；乘法顺序错了就消不掉。}
\end{solutionblock}

\begin{problemblock}
\textbf{例 19.} 设 \(A,B,C\) 均为 \(n\) 阶矩阵。若 \(AB=C\)，且 \(B\) 可逆，则（\quad）：
\[
\begin{array}{l}
\text{A. 矩阵 }C\text{ 的行向量组与矩阵 }A\text{ 的行向量组等价},\\
\text{B. 矩阵 }C\text{ 的列向量组与矩阵 }A\text{ 的列向量组等价},\\
\text{C. 矩阵 }C\text{ 的行向量组与矩阵 }B\text{ 的行向量组等价},\\
\text{D. 矩阵 }C\text{ 的列向量组与矩阵 }B\text{ 的列向量组等价}.
\end{array}
\]
\end{problemblock}
\begin{solutionblock}
\analysis{右乘可逆矩阵 \(B\) 等价于对 \(A\) 的列向量组作可逆线性组合。也就是说 \(C=AB\) 的每一列都是 \(A\) 的列向量线性组合；反过来 \(A=CB^{-1}\)，\(A\) 的每一列也是 \(C\) 的列向量线性组合。因此 \(A,C\) 的列向量组等价。}
\answer{B}
\examnote{右乘可逆矩阵改变列向量组的表示但不改变列空间；左乘可逆矩阵才对应行向量组等价。}
\end{solutionblock}

\begin{problemblock}
\textbf{例 20.} 设
\[
A=(\alpha_1,\alpha_2,\alpha_3),\qquad B=(\beta_1,\beta_2,\beta_3)
\]
都是三阶矩阵，规定三阶矩阵
\[
C=\begin{pmatrix}
\alpha_1^T\beta_1&\alpha_1^T\beta_2&\alpha_1^T\beta_3\\
\alpha_2^T\beta_1&\alpha_2^T\beta_2&\alpha_2^T\beta_3\\
\alpha_3^T\beta_1&\alpha_3^T\beta_2&\alpha_3^T\beta_3
\end{pmatrix}.
\]
证明：\(C\) 可逆的充分必要条件是 \(A,B\) 都可逆。
\end{problemblock}
\begin{solutionblock}
\analysis{注意 \(A\) 的列为 \(\alpha_i\)，所以
\[
A^T=\begin{pmatrix}
\alpha_1^T\\
\alpha_2^T\\
\alpha_3^T
\end{pmatrix}.
\]
于是题中矩阵正是
\[
C=A^TB.
\]
取行列式得
\[
|C|=|A^T|\,|B|=|A|\,|B|.
\]
因此 \(C\) 可逆等价于 \(|C|\ne0\)，等价于
\[
|A|\ne0,\qquad |B|\ne0,
\]
即 \(A,B\) 都可逆。反向也由同一等式立即成立。}
\answer{\(C\) 可逆 \(\Longleftrightarrow A,B\) 都可逆}
\examnote{这类由内积排成的矩阵常写成 \(A^TB\)。先识别矩阵乘法，比逐个元素讨论高效。}
\end{solutionblock}

\begin{problemblock}
\textbf{例 21.} 设
\[
A=\begin{pmatrix}
1&0&0&0\\
-2&3&0&0\\
0&-4&5&0\\
0&0&-6&7
\end{pmatrix},
\]
\(E\) 为四阶单位矩阵，且
\[
B=(E+A)^{-1}(E-A).
\]
求 \((E+B)^{-1}\)。
\end{problemblock}
\begin{solutionblock}
\analysis{由定义
\[
E+B=E+(E+A)^{-1}(E-A).
\]
把 \(E\) 写成 \((E+A)^{-1}(E+A)\)，则
\[
E+B=(E+A)^{-1}\bigl[(E+A)+(E-A)\bigr]
=2(E+A)^{-1}.
\]
所以
\[
(E+B)^{-1}=\frac12(E+A).
\]
代入题中 \(A\)：
\[
\frac12(E+A)=
\begin{pmatrix}
1&0&0&0\\
-1&2&0&0\\
0&-2&3&0\\
0&0&-3&4
\end{pmatrix}.
\]}
\answer{\(\displaystyle
\begin{pmatrix}
1&0&0&0\\
-1&2&0&0\\
0&-2&3&0\\
0&0&-3&4
\end{pmatrix}\)}
\examnote{形如 \((E+A)^{-1}(E-A)\) 的 Cayley 变换题，先整体代数化简，不要先求逆矩阵。}
\end{solutionblock}

\begin{problemblock}
\textbf{例 22.} 设 \(\alpha,\beta\) 是相互正交的 \(n\) 维列向量，\(E\) 是 \(n\) 阶单位矩阵，
\[
A=E+\alpha\beta^T.
\]
求 \(A^{-1}\)。
\end{problemblock}
\begin{solutionblock}
\analysis{相互正交说明
\[
\beta^T\alpha=0.
\]
尝试
\[
(E+\alpha\beta^T)(E-\alpha\beta^T)
=E-\alpha\beta^T+\alpha\beta^T-\alpha(\beta^T\alpha)\beta^T
=E.
\]
同理反向乘积也为 \(E\)，故
\[
A^{-1}=E-\alpha\beta^T.
\]}
\answer{\(E-\alpha\beta^T\)}
\examnote{秩一扰动 \(E+uv^T\) 的逆常用 Sherman--Morrison 公式；这里因 \(v^Tu=0\)，公式极简。}
\end{solutionblock}

\begin{problemblock}
\textbf{例 23.} 设 \(n\) 阶行列式
\[
|A|=
\begin{vmatrix}
0&0&\cdots&0&-1\\
-1&0&\cdots&0&0\\
0&-1&\cdots&0&0\\
\vdots&\vdots&\ddots&\vdots&\vdots\\
0&0&\cdots&-1&0
\end{vmatrix},\qquad n>2,
\]
\(A_{ij}\) 为元素 \(a_{ij}\) 的代数余子式，则
\[
\sum_{i=1}^n\sum_{j=1}^n A_{ij}=(\quad)
\]
\[
\text{A. }n,\qquad \text{B. }-n,\qquad \text{C. }(-1)^n n,\qquad \text{D. }(-1)^{n-1}n.
\]
\end{problemblock}
\begin{solutionblock}
\analysis{矩阵 \(A\) 是带负号的循环置换矩阵，每一行每一列只有一个 \(-1\)。它可逆，且 \(A^{-1}\mathbf1=-\mathbf1\)，其中 \(\mathbf1=(1,\ldots,1)^T\)。又
\[
(A_{ij})^T=A^*=|A|A^{-1}.
\]
因此所有代数余子式之和为
\[
\sum_{i,j}A_{ij}
=\mathbf1^T(A_{ij})\mathbf1
=\mathbf1^T(A^*)^T\mathbf1
=|A|\,\mathbf1^T(A^{-1})^T\mathbf1.
\]
该循环置换的行列式为 \(|A|=(-1)^n\cdot(-1)^{n-1}=-1\)，而
\[
(A^{-1})^T\mathbf1=-\mathbf1.
\]
所以
\[
\sum_{i,j}A_{ij}=(-1)\cdot(-n)=n.
\]}
\method{方法二：扰动法}
\(\sum A_{ij}\) 是 \(|A+tJ|\) 在 \(t=0\) 处对 \(t\) 的导数，其中 \(J=\mathbf1\mathbf1^T\)。用行列式引理可得同样结果 \(n\)。
\answer{A；\(n\)}
\examnote{代数余子式求和常转成 \(\mathbf1^TA^*\mathbf1\)，这是伴随矩阵的典型用法。}
\end{solutionblock}

\begin{problemblock}
\textbf{例 24.} 设 \(E\) 为三阶单位矩阵，\(A=(a_{ij})\) 为三阶非零实矩阵，\(A_{ij}\) 是 \(|A|\) 中元素 \(a_{ij}\) 的代数余子式。若
\[
A_{ij}=a_{ij}\quad(i,j=1,2,3),
\]
则（\quad）：
\[
\text{A. }A-E\text{ 不可逆},\quad
\text{B. }A-E\text{ 可逆},\quad
\text{C. }A+E\text{ 不可逆},\quad
\text{D. }A+E\text{ 可逆}.
\]
\end{problemblock}
\begin{solutionblock}
\analysis{设余子式矩阵为 \(C=(A_{ij})\)，题设给出 \(C=A\)。伴随矩阵为
\[
A^*=C^T=A^T.
\]
由伴随矩阵恒等式
\[
AA^*=|A|E
\]
得
\[
AA^T=|A|E.
\]
若 \(|A|=0\)，则 \(AA^T=O\)，从而 \(A=O\)，与 \(A\) 非零矛盾。因此 \(|A|\ne0\)。两边取行列式：
\[
|A|^2=|A|^3,
\]
故 \(|A|=1\)。于是
\[
AA^T=E,
\]
即 \(A\) 是三阶正交矩阵，且 \(|A|=1\)。三阶实正交矩阵的特征值中，非实特征值必成共轭对，乘积为 \(1\)，又总行列式为 \(1\)，所以必有实特征值 \(1\)。因此
\[
|A-E|=0,
\]
即 \(A-E\) 不可逆。}
\answer{A}
\examnote{三阶正交矩阵且行列式为 \(1\) 时，必有特征值 \(1\)。这是本题从伴随矩阵走向特征值的关键。}
\end{solutionblock}

\begin{problemblock}
\textbf{例 25.} 设 \(A,B\) 为 \(n\) 阶可逆矩阵，\(E\) 为 \(n\) 阶单位矩阵，\(M^*\) 为矩阵 \(M\) 的伴随矩阵。求
\[
\begin{pmatrix}A&E\\O&B\end{pmatrix}^*.
\]
\end{problemblock}
\begin{solutionblock}
\analysis{记
\[
M=\begin{pmatrix}A&E\\O&B\end{pmatrix}.
\]
它是分块上三角矩阵，
\[
|M|=|A|\,|B|.
\]
又
\[
M^{-1}=
\begin{pmatrix}
A^{-1}&-A^{-1}B^{-1}\\
O&B^{-1}
\end{pmatrix}.
\]
由 \(M^*=|M|M^{-1}\)，并用
\[
A^*=|A|A^{-1},\qquad B^*=|B|B^{-1},
\]
得到
\[
M^*=
\begin{pmatrix}
|B|A^*&-A^*B^*\\
O&|A|B^*
\end{pmatrix}.
\]}
\answer{\(\displaystyle
\begin{pmatrix}
|B|A^*&-A^*B^*\\
O&|A|B^*
\end{pmatrix}\)}
\examnote{可逆矩阵的伴随矩阵直接用 \(M^*=|M|M^{-1}\)，分块上三角矩阵先求逆最省力。}
\end{solutionblock}

\begin{problemblock}
\textbf{例 26.} 设 \(n\) 阶矩阵 \(A\) 为反对称矩阵，\(r(A)=n\)，\(\beta\) 为非零列向量。证明：
\[
\begin{pmatrix}A&\beta\\ \beta^T&0\end{pmatrix}x=0
\]
必有非零解。
\end{problemblock}
\begin{solutionblock}
\analysis{因为 \(r(A)=n\)，所以 \(A\) 可逆。反对称矩阵满足 \(A^T=-A\)，于是
\[
(A^{-1})^T=(A^T)^{-1}=(-A)^{-1}=-A^{-1},
\]
即 \(A^{-1}\) 仍为反对称矩阵。对任意列向量 \(\beta\)，标量
\[
\beta^TA^{-1}\beta
\]
等于自身转置：
\[
(\beta^TA^{-1}\beta)^T=\beta^T(A^{-1})^T\beta=-\beta^TA^{-1}\beta.
\]
故
\[
\beta^TA^{-1}\beta=0.
\]
用 Schur 补计算分块矩阵行列式：
\[
\left|\begin{matrix}A&\beta\\\beta^T&0\end{matrix}\right|
 =
|A|\left(0-\beta^TA^{-1}\beta\right)=0.
\]
系数矩阵奇异，因此齐次线性方程组必有非零解。}
\answer{已证明存在非零解}
\examnote{反对称矩阵的核心性质是 \(x^TSx=0\)。分块矩阵是否奇异，用 Schur 补一行解决。}
\end{solutionblock}

\begin{problemblock}
\textbf{例 27.} 设 \(A,B\) 均为 \(n\) 阶正交矩阵，且
\[
|A|+|B|=0.
\]
证明：
\[
|A+B|=0.
\]
\end{problemblock}
\begin{solutionblock}
\analysis{因为 \(A\) 正交，所以 \(A^{-1}=A^T\)，且 \(|A|=\pm1\)。由 \(|A|+|B|=0\) 得 \(|B|=-|A|\)。写成
\[
A+B=A(E+A^TB).
\]
令
\[
Q=A^TB.
\]
则 \(Q\) 仍是正交矩阵，且
\[
|Q|=|A^T||B|=|A|\,|B|=-1.
\]
实正交矩阵的特征值模长均为 \(1\)，非实特征值成共轭对，乘积为 \(1\)。若 \(|Q|=-1\)，则 \(Q\) 必有特征值 \(-1\)。因此 \(E+Q\) 不可逆，
\[
|E+Q|=0.
\]
于是
\[
|A+B|=|A|\,|E+Q|=0.
\]}
\answer{\(|A+B|=0\)}
\examnote{两个正交矩阵相加，常先左乘一个正交矩阵的逆，化为 \(E+Q\)。}
\end{solutionblock}

\begin{problemblock}
\textbf{例 28.} 设 \(\alpha,\beta\) 为三维列向量，\(\beta^T\) 为 \(\beta\) 的转置。若矩阵 \(\alpha\beta^T\) 相似于
\[
\begin{pmatrix}2&0&0\\0&0&0\\0&0&0\end{pmatrix},
\]
求 \(\beta^T\alpha\)。
\end{problemblock}
\begin{solutionblock}
\analysis{相似矩阵有相同迹。秩一矩阵 \(\alpha\beta^T\) 的迹为
\[
\operatorname{tr}(\alpha\beta^T)=\beta^T\alpha.
\]
而题给相似标准矩阵的迹为
\[
2+0+0=2.
\]
所以
\[
\beta^T\alpha=2.
\]}
\answer{\(2\)}
\examnote{秩一矩阵 \(\alpha\beta^T\) 的非零特征值、迹都等于 \(\beta^T\alpha\)。}
\end{solutionblock}

\begin{problemblock}
\textbf{例 29.} 已知
\[
\alpha=(1,2,3),\qquad \beta=\left(1,\frac12,\frac13\right),
\]
设
\[
A=\alpha^T\beta,
\]
其中 \(\alpha^T\) 是 \(\alpha\) 的转置。求 \(A^n\)。
\end{problemblock}
\begin{solutionblock}
\analysis{这里 \(A=\alpha^T\beta\) 是 \(3\times3\) 秩一矩阵。利用秩一矩阵乘法：
\[
A^2=(\alpha^T\beta)(\alpha^T\beta)
=\alpha^T(\beta\alpha^T)\beta.
\]
其中
\[
\beta\alpha^T=1\cdot1+\frac12\cdot2+\frac13\cdot3=3.
\]
所以
\[
A^2=3A.
\]
递推得到
\[
A^n=3^{n-1}A\qquad(n\ge1).
\]
即
\[
A^n=3^{n-1}\alpha^T\beta.
\]}
\answer{\(A^n=3^{\,n-1}A=3^{\,n-1}\alpha^T\beta\quad(n\ge1)\)}
\examnote{秩一矩阵 \(uv^T\) 满足 \((uv^T)^n=(v^Tu)^{n-1}uv^T\)。}
\end{solutionblock}

\begin{problemblock}
\textbf{例 30.} 求二次型
\[
f(x_1,x_2,x_3)=\sum_{i=1}^3\sum_{j=1}^3 ijx_ix_j
\]
所对应的二次型矩阵的特征值及特征向量。
\end{problemblock}
\begin{solutionblock}
\analysis{二次型矩阵为
\[
A=(ij)_{3\times3}
=\begin{pmatrix}
1&2&3\\
2&4&6\\
3&6&9
\end{pmatrix}
=vv^T,\qquad v=(1,2,3)^T.
\]
这是秩一实对称矩阵。非零特征值为
\[
v^Tv=1^2+2^2+3^2=14,
\]
对应特征向量为 \(v=(1,2,3)^T\)。其余特征值为 \(0,0\)，对应特征子空间为
\[
v^Tx=x_1+2x_2+3x_3=0.
\]
可取一组线性无关特征向量
\[
(-2,1,0)^T,\qquad (-3,0,1)^T.
\]}
\answer{特征值 \(14,0,0\)。\(\lambda=14\) 的特征向量可取 \((1,2,3)^T\)；\(\lambda=0\) 的特征向量满足 \(x_1+2x_2+3x_3=0\)，如 \((-2,1,0)^T,(-3,0,1)^T\)。}
\examnote{矩阵形如 \(vv^T\) 时，直接用秩一矩阵结论，不必展开三阶特征行列式。}
\end{solutionblock}

\begin{problemblock}
\textbf{例 31.} 设
\[
A=\begin{pmatrix}
1&-1&-1&-1\\
-1&1&-1&-1\\
-1&-1&1&-1\\
-1&-1&-1&1
\end{pmatrix},
\]
求
\[
\sum_{i=1}^4\sum_{j=1}^4 A_{ij},
\]
其中 \(A_{ij}\) 为元素 \(a_{ij}\) 的代数余子式。
\end{problemblock}
\begin{solutionblock}
\analysis{矩阵可写成
\[
A=2E-J,
\]
其中 \(J=\mathbf1\mathbf1^T\)，\(\mathbf1=(1,1,1,1)^T\)。在 \(\mathbf1\) 方向上，
\[
A\mathbf1=(2E-J)\mathbf1=(2-4)\mathbf1=-2\mathbf1;
\]
在与 \(\mathbf1\) 正交的三维子空间上，\(Jx=0\)，所以特征值为 \(2\)。因此
\[
|A|=(-2)\cdot2^3=-16.
\]
所有代数余子式之和为
\[
\mathbf1^T(A^*)^T\mathbf1.
\]
由于 \(A\) 对称，\(A^*=|A|A^{-1}\) 也对称，故
\[
\sum_{i,j}A_{ij}
=\mathbf1^TA^*\mathbf1
=|A|\mathbf1^TA^{-1}\mathbf1.
\]
又 \(A^{-1}\mathbf1=-\dfrac12\mathbf1\)，所以
\[
\mathbf1^TA^{-1}\mathbf1=-\frac12\cdot4=-2.
\]
于是
\[
\sum_{i,j}A_{ij}=(-16)(-2)=32.
\]}
\answer{\(32\)}
\examnote{\(a\) 在对角线、\(b\) 在非对角线的矩阵，优先写成 \((a-b)E+bJ\)。}
\end{solutionblock}

\begin{problemblock}
\textbf{例 32.} 计算
\[
A=
\begin{pmatrix}1&0&0\\0&0&1\\0&1&0\end{pmatrix}^{9}
\begin{pmatrix}1&2&3\\4&5&6\\7&8&9\end{pmatrix}
\begin{pmatrix}1&0&0\\0&2&0\\0&0&1\end{pmatrix}^{10}.
\]
\end{problemblock}
\begin{solutionblock}
\analysis{左侧矩阵
\[
P=\begin{pmatrix}1&0&0\\0&0&1\\0&1&0\end{pmatrix}
\]
满足 \(P^2=E\)，故 \(P^9=P\)，作用是交换第 2、3 行。右侧对角矩阵
\[
D=\operatorname{diag}(1,2,1)
\]
满足
\[
D^{10}=\operatorname{diag}(1,2^{10},1)=\operatorname{diag}(1,1024,1),
\]
作用是把第 2 列乘以 \(1024\)。先右乘：
\[
\begin{pmatrix}1&2&3\\4&5&6\\7&8&9\end{pmatrix}
D^{10}
=
\begin{pmatrix}
1&2048&3\\
4&5120&6\\
7&8192&9
\end{pmatrix}.
\]
再左乘 \(P\)，交换第 2、3 行，得
\[
A=
\begin{pmatrix}
1&2048&3\\
7&8192&9\\
4&5120&6
\end{pmatrix}.
\]}
\answer{\(\displaystyle
\begin{pmatrix}
1&2048&3\\
7&8192&9\\
4&5120&6
\end{pmatrix}\)}
\examnote{置换矩阵左乘换行，右乘换列；对角矩阵右乘按列缩放。}
\end{solutionblock}

\begin{problemblock}
\textbf{例 33.} 设 \(A\) 为 \(n\ge2\) 阶可逆矩阵，交换 \(A\) 的第 1 行与第 2 行得到矩阵 \(B\)。\(A^*,B^*\) 分别为 \(A,B\) 的伴随矩阵，则（\quad）：
\[
\begin{array}{ll}
\text{A. 交换 }A^*\text{ 的第 1 列与第 2 列得 }B^*,&
\text{B. 交换 }A^*\text{ 的第 1 行与第 2 行得 }B^*,\\
\text{C. 交换 }A^*\text{ 的第 1 列与第 2 列得 }-B^*,&
\text{D. 交换 }A^*\text{ 的第 1 行与第 2 行得 }-B^*.
\end{array}
\]
\end{problemblock}
\begin{solutionblock}
\analysis{设 \(P\) 是交换第 1、2 行的初等置换矩阵，则
\[
B=PA,\qquad |P|=-1,\qquad P^{-1}=P.
\]
因为 \(A,B\) 可逆，
\[
B^*=|B|B^{-1}=|P||A|A^{-1}P^{-1}=-A^*P.
\]
右乘 \(P\) 表示交换矩阵的第 1、2 列，所以
\[
A^*P
\]
正是交换 \(A^*\) 的第 1 列与第 2 列所得矩阵。由上式
\[
A^*P=-B^*.
\]
故交换 \(A^*\) 的第 1 列与第 2 列得 \(-B^*\)。}
\answer{C}
\examnote{原矩阵换行，对伴随矩阵表现为换列并带行列式符号。可用 \(B=PA\) 一步推出。}
\end{solutionblock}

\begin{problemblock}
\textbf{例 34.} 求
\[
\begin{pmatrix}
0&1&0\\
-1&0&1\\
0&-1&0
\end{pmatrix}^{2017}.
\]
\end{problemblock}
\begin{solutionblock}
\analysis{记
\[
S=\begin{pmatrix}
0&1&0\\
-1&0&1\\
0&-1&0
\end{pmatrix}.
\]
直接计算可得
\[
S^3=-2S.
\]
于是对奇数次幂有递推：
\[
S^{2m+1}=(-2)^mS\qquad(m\ge0).
\]
因为
\[
2017=2\cdot1008+1,
\]
所以
\[
S^{2017}=(-2)^{1008}S=2^{1008}S.
\]}
\answer{\(\displaystyle
2^{1008}\begin{pmatrix}
0&1&0\\
-1&0&1\\
0&-1&0
\end{pmatrix}\)}
\examnote{高次矩阵幂先找低次递推关系，如 \(S^3=cS\)，再降幂。}
\end{solutionblock}

\begin{problemblock}
\textbf{例 35.} 设
\[
A=\begin{pmatrix}
1&0&0\\
1&0&1\\
0&1&0
\end{pmatrix}.
\]
\begin{enumerate}
\item 证明当 \(n>1\) 时，\(A^n=A^{n-2}+A^2-E\)；
\item 求 \(A^n\)。
\end{enumerate}
\end{problemblock}
\begin{solutionblock}
\analysis{先计算
\[
A^2=\begin{pmatrix}
1&0&0\\
1&1&0\\
1&0&1
\end{pmatrix}.
\]
继续计算
\[
A^3=
\begin{pmatrix}
1&0&0\\
2&0&1\\
1&1&0
\end{pmatrix}.
\]
由特征多项式
\[
|\lambda E-A|=(\lambda-1)^2(\lambda+1)
\]
得
\[
A^3-A^2-A+E=O.
\]
两边右乘 \(A^{n-3}\)（\(n>1\) 时可按 \(n=2\) 直接验证，其余由递推）可得
\[
A^n=A^{n-2}+A^2-E.
\]
下面求通项。由前几项观察并归纳：
\[
A^{2k}=
\begin{pmatrix}
1&0&0\\
k&1&0\\
k&0&1
\end{pmatrix}\quad(k\ge1),
\]
\[
A^{2k+1}=
\begin{pmatrix}
1&0&0\\
k+1&0&1\\
k&1&0
\end{pmatrix}\quad(k\ge0).
\]
把它们代入递推式 \(A^n=A^{n-2}+A^2-E\) 可完成归纳证明。}
\method{方法二：分解为置换加幂零}
矩阵在后两个坐标上有交换结构，偶次幂回到对角结构，奇次幂保留交换结构；再追踪第一列累加量即可得到同样的偶奇通项。
\answer{\[
A^{2k}=
\begin{pmatrix}
1&0&0\\
k&1&0\\
k&0&1
\end{pmatrix}\ (k\ge1),\qquad
A^{2k+1}=
\begin{pmatrix}
1&0&0\\
k+1&0&1\\
k&1&0
\end{pmatrix}\ (k\ge0).
\]}
\examnote{矩阵高次幂常先用 Cayley--Hamilton 建递推，再按奇偶或特征根结构求通项。}
\end{solutionblock}

\begin{problemblock}
\textbf{例 36.} 设
\[
B=\begin{pmatrix}
0&1&0&0\\
0&0&1&0\\
0&0&0&1\\
0&0&0&0
\end{pmatrix},\qquad
A=E+B+B^2+B^3.
\]
求 \(A^{-1}\)。
\end{problemblock}
\begin{solutionblock}
\analysis{矩阵 \(B\) 是四阶幂零 Jordan 块，满足
\[
B^4=O.
\]
因此
\[
(E-B)(E+B+B^2+B^3)=E-B^4=E.
\]
所以
\[
A^{-1}=E-B.
\]
写成矩阵为
\[
A^{-1}=
\begin{pmatrix}
1&-1&0&0\\
0&1&-1&0\\
0&0&1&-1\\
0&0&0&1
\end{pmatrix}.
\]}
\answer{\(\displaystyle
\begin{pmatrix}
1&-1&0&0\\
0&1&-1&0\\
0&0&1&-1\\
0&0&0&1
\end{pmatrix}\)}
\examnote{有限等比和 \(E+B+\cdots+B^{m}\) 的逆通常由 \((E-B)\) 给出，只要 \(B^{m+1}=O\)。}
\end{solutionblock}

\begin{problemblock}
\textbf{例 37.} 设
\[
A=\begin{pmatrix}
1&2&3\\
2&4&6\\
3&6&9
\end{pmatrix},
\]
求 \(A^{2027}\)。
\end{problemblock}
\begin{solutionblock}
\analysis{矩阵
\[
A=vv^T,\qquad v=(1,2,3)^T.
\]
因此
\[
A^2=vv^Tvv^T=(v^Tv)vv^T.
\]
而
\[
v^Tv=1^2+2^2+3^2=14.
\]
所以
\[
A^2=14A.
\]
递推得
\[
A^m=14^{m-1}A\qquad(m\ge1).
\]
取 \(m=2027\)，
\[
A^{2027}=14^{2026}A.
\]}
\answer{\(\displaystyle
14^{2026}\begin{pmatrix}
1&2&3\\
2&4&6\\
3&6&9
\end{pmatrix}\)}
\examnote{秩一对称矩阵 \(vv^T\) 的高次幂用 \(v^Tv\) 递推。}
\end{solutionblock}

\begin{problemblock}
\textbf{例 38.} 设
\[
A=\begin{pmatrix}
2&1&1\\
2&3&2\\
3&3&4
\end{pmatrix}.
\]
求：
\[
(1)\ A^n;\qquad (2)\ A^{-1};\qquad (3)\ A\text{ 的特征值和特征向量}.
\]
\end{problemblock}
\begin{solutionblock}
\analysis{注意
\[
A=E+uv^T,\qquad
u=(1,2,3)^T,\quad v=(1,1,1)^T.
\]
且
\[
v^Tu=1+2+3=6.
\]
对秩一扰动，有
\[
(E+uv^T)^n
=E+\frac{(1+v^Tu)^n-1}{v^Tu}\,uv^T.
\]
因此
\[
A^n
=E+\frac{7^n-1}{6}
\begin{pmatrix}
1&1&1\\
2&2&2\\
3&3&3
\end{pmatrix}.
\]
逆矩阵由 Sherman--Morrison 公式：
\[
A^{-1}=E-\frac1{1+v^Tu}uv^T
=E-\frac17
\begin{pmatrix}
1&1&1\\
2&2&2\\
3&3&3
\end{pmatrix},
\]
即
\[
A^{-1}=
\begin{pmatrix}
\frac67&-\frac17&-\frac17\\
-\frac27&\frac57&-\frac27\\
-\frac37&-\frac37&\frac47
\end{pmatrix}.
\]
特征值方面，\(u\) 方向上
\[
Au=(E+uv^T)u=(1+v^Tu)u=7u,
\]
故 \(\lambda=7\) 的特征向量可取 \((1,2,3)^T\)。若 \(v^Tx=x_1+x_2+x_3=0\)，则 \(uv^Tx=0\)，所以 \(Ax=x\)。因此 \(\lambda=1\) 的特征子空间为
\[
x_1+x_2+x_3=0,
\]
可取基
\[
(-1,1,0)^T,\qquad (-1,0,1)^T.
\]}
\method{方法二：对角化}
因为 \(\lambda=1\) 有两个线性无关特征向量，\(\lambda=7\) 有一个特征向量，三阶矩阵可对角化。取
\[
P=\begin{pmatrix}
1&-1&-1\\
2&1&0\\
3&0&1
\end{pmatrix},
\]
则 \(P^{-1}AP=\operatorname{diag}(7,1,1)\)，从而 \(A^n=P\operatorname{diag}(7^n,1,1)P^{-1}\)，化简即上式。
\answer{\[
A^n=E+\frac{7^n-1}{6}
\begin{pmatrix}
1&1&1\\
2&2&2\\
3&3&3
\end{pmatrix},
\quad
A^{-1}=
\begin{pmatrix}
\frac67&-\frac17&-\frac17\\
-\frac27&\frac57&-\frac27\\
-\frac37&-\frac37&\frac47
\end{pmatrix}.
\]
特征值为 \(7,1,1\)。\(\lambda=7\) 的特征向量可取 \((1,2,3)^T\)；\(\lambda=1\) 的特征向量满足 \(x_1+x_2+x_3=0\)。}
\examnote{看到 \(A-E\) 每一列相同或每一行成比例，应立即想到秩一扰动。}
\end{solutionblock}

\begin{problemblock}
\textbf{例 39.} 已知
\[
A=\begin{pmatrix}
0&1&0&0\\
1&0&0&0\\
0&0&1&-1\\
0&0&-1&1
\end{pmatrix},
\]
求 \(A^5\)。
\end{problemblock}
\begin{solutionblock}
\analysis{矩阵为分块对角形：
\[
A=\begin{pmatrix}
P&O\\
O&C
\end{pmatrix},
\qquad
P=\begin{pmatrix}0&1\\1&0\end{pmatrix},
\quad
C=\begin{pmatrix}1&-1\\-1&1\end{pmatrix}.
\]
其中 \(P^2=E\)，故 \(P^5=P\)。矩阵 \(C\) 的特征值为 \(0,2\)，且
\[
C^2=2C,
\]
故
\[
C^5=2^4C=16C.
\]
于是
\[
A^5=
\begin{pmatrix}
0&1&0&0\\
1&0&0&0\\
0&0&16&-16\\
0&0&-16&16
\end{pmatrix}.
\]}
\answer{\(\displaystyle
\begin{pmatrix}
0&1&0&0\\
1&0&0&0\\
0&0&16&-16\\
0&0&-16&16
\end{pmatrix}\)}
\examnote{分块对角矩阵的幂等于各块分别取幂。}
\end{solutionblock}

\begin{problemblock}
\textbf{例 40.} 设 \(n\) 阶矩阵 \(A,B\) 满足
\[
AB=aA+bB,\qquad ab\ne0.
\]
证明：
\[
(1)\ A-bE\text{ 和 }B-aE\text{ 都可逆};\qquad
(2)\ AB=BA.
\]
\end{problemblock}
\begin{solutionblock}
\analysis{由题设
\[
AB-aA-bB=O.
\]
于是
\[
(A-bE)(B-aE)=AB-aA-bB+abE=abE.
\]
因为 \(ab\ne0\)，所以 \((A-bE)(B-aE)\) 是非零数量倍单位矩阵。对方阵而言，两个方阵乘积可逆则每个因子均可逆，因此 \(A-bE\) 与 \(B-aE\) 都可逆。

并且由
\[
(A-bE)(B-aE)=abE
\]
可知
\[
B-aE=ab(A-bE)^{-1}.
\]
所以 \(A-bE\) 与 \(B-aE\) 可交换，进而
\[
(B-aE)(A-bE)=abE.
\]
展开得
\[
BA-bB-aA+abE=abE,
\]
故
\[
BA=aA+bB=AB.
\]}
\answer{已证明 \(A-bE,B-aE\) 均可逆，且 \(AB=BA\)}
\examnote{把 \(AB=aA+bB\) 移项并配成 \((A-bE)(B-aE)\)，是本题唯一关键。}
\end{solutionblock}

\begin{problemblock}
\textbf{例 41.} 设 \(A,B,C\) 均为 \(n\) 阶矩阵，且
\[
AB=O,\qquad A+BC=E.
\]
则（\quad）。
\[
\begin{array}{ll}
\text{A. }r(A)=n, & \text{B. }r(B)=n,\\
\text{C. }r(A)+r(B)<n, & \text{D. }r(A)+r(B)<n.
\end{array}
\]
\end{problemblock}
\begin{solutionblock}
\analysis{由
\[
A+BC=E
\]
两边左乘 \(A\)，得
\[
A^2+ABC=A.
\]
而 \(AB=O\)，所以 \(ABC=O\)，故
\[
A^2=A.
\]
因此 \(A\) 是幂等矩阵，\(E-A=BC\)，并且
\[
r(E-A)=n-r(A).
\]
另一方面
\[
AB=O
\]
说明 \(B\) 的列空间包含在 \(A\) 的零空间中，所以
\[
r(B)\le n-r(A).
\]
又
\[
n-r(A)=r(E-A)=r(BC)\le r(B).
\]
两边合并得到
\[
r(B)=n-r(A),
\]
即
\[
r(A)+r(B)=n.
\]
因此精确结论是 \(r(A)+r(B)=n\)。若原题选项 C、D 均印为 \(r(A)+r(B)<n\)，则按题设无正确选项，疑似应有一个选项为等号。}
\answer{\(r(A)+r(B)=n\)。若按 PDF 中 C、D 均为 \(<n\) 的印刷，则四个选项无正确项。}
\examnote{这题不是严格小于，而是等于。用 \(E-A=BC\) 与 \(AB=O\) 从上下界夹出 \(r(B)=n-r(A)\)。}
\end{solutionblock}

\begin{problemblock}
\textbf{例 42.} 设 \(A\) 为 \(m\times n\) 矩阵，\(B\) 为 \(n\times m\) 矩阵，\(E\) 为 \(m\) 阶单位矩阵。若
\[
AB=E,
\]
则（\quad）：
\[
\begin{array}{ll}
\text{A. }r(A)=m,\ r(B)=m, & \text{B. }r(A)=m,\ r(B)=n,\\
\text{C. }r(A)=n,\ r(B)=m, & \text{D. }r(A)=n,\ r(B)=n.
\end{array}
\]
\end{problemblock}
\begin{solutionblock}
\analysis{由 \(AB=E_m\)，得
\[
r(AB)=r(E_m)=m.
\]
又
\[
r(AB)\le r(A)\le m,
\]
所以
\[
r(A)=m.
\]
同理
\[
r(AB)\le r(B)\le m,
\]
所以
\[
r(B)=m.
\]
因此选 A。顺便可知 \(m\le n\)，否则 \(m\times n\) 矩阵 \(A\) 不可能有行满秩 \(m\)。}
\answer{A；\(r(A)=m,\ r(B)=m\)}
\examnote{矩阵乘积等于单位矩阵时，两个因子都必须达到相应的满秩。}
\end{solutionblock}

\begin{problemblock}
\textbf{例 43.} \(A\) 为四阶矩阵，若
\[
A(A-A^*)=O,\qquad A\ne A^*,
\]
则 \(r(A)\) 可能为（\quad）：
\[
\text{A. }0\text{ 或 }1,\qquad
\text{B. }1\text{ 或 }3,\qquad
\text{C. }2\text{ 或 }3,\qquad
\text{D. }1\text{ 或 }2.
\]
\end{problemblock}
\begin{solutionblock}
\analysis{分秩讨论。若 \(r(A)=4\)，则 \(A\) 可逆。由
\[
A(A-A^*)=O
\]
左乘 \(A^{-1}\) 得
\[
A=A^*,
\]
与题设矛盾，所以 \(r(A)\ne4\)。

若 \(r(A)=3\)，则 \(|A|=0\)，但 \(A^*\ne O\) 且 \(AA^*=O\)。原式化为
\[
A^2=AA^*=O.
\]
然而 \(A^2=O\) 时 \(\operatorname{Im}A\subset\operatorname{Ker}A\)，故
\[
r(A)\le 4-r(A),
\]
即 \(r(A)\le2\)，与 \(r(A)=3\) 矛盾。

若 \(r(A)\le2\)，四阶矩阵的伴随矩阵满足 \(A^*=O\)，原式化为
\[
A^2=O.
\]
非零平方零矩阵的秩可以是 \(1\) 或 \(2\)。又 \(A\ne A^*=O\)，所以 \(A\ne O\)，排除 \(r(A)=0\)。因此可能的秩为 \(1\) 或 \(2\)。}
\answer{D；\(1\) 或 \(2\)}
\examnote{四阶矩阵中，\(r(A)\le2\Rightarrow A^*=O\)，\(r(A)=3\Rightarrow r(A^*)=1\)。伴随矩阵的秩分类很常考。}
\end{solutionblock}

\begin{problemblock}
\textbf{例 44.} 设 \(n\) 阶矩阵 \(A,B,C\) 满足
\[
r(A)+r(B)+r(C)=r(ABC)+2n.
\]
给出下列四个结论：
\[
\begin{array}{ll}
\text{(1) }r(ABC)+n=r(AB)+r(C),&
\text{(2) }r(AB)+n=r(A)+r(B),\\
\text{(3) }r(A)=r(B)=r(C)=n,&
\text{(4) }r(AB)=r(BC)=n.
\end{array}
\]
其中正确的选项是（\quad）：
\[
\text{A. (1)(2)},\qquad \text{B. (1)(3)},\qquad \text{C. (2)(4)},\qquad \text{D. (3)(4)}.
\]
\end{problemblock}
\begin{solutionblock}
\analysis{使用 Frobenius 秩不等式：
\[
r(AB)\ge r(A)+r(B)-n,
\]
\[
r(ABC)\ge r(AB)+r(C)-n.
\]
两式相加的链式下界给出
\[
r(ABC)\ge r(A)+r(B)+r(C)-2n.
\]
题设正好是
\[
r(ABC)=r(A)+r(B)+r(C)-2n,
\]
说明上述两步不等式都取等号。因此
\[
r(AB)=r(A)+r(B)-n,
\]
即
\[
r(AB)+n=r(A)+r(B),
\]
结论 (2) 正确；并且
\[
r(ABC)=r(AB)+r(C)-n,
\]
即
\[
r(ABC)+n=r(AB)+r(C),
\]
结论 (1) 正确。题设并不强迫 \(A,B,C\) 都满秩，也不强迫 \(AB,BC\) 满秩，因此 (3)(4) 不一定正确。}
\answer{A；(1)(2)}
\examnote{等号型秩题常是 Frobenius 不等式连续取等，先写不等式链，再看题设正好达到下界。}
\end{solutionblock}

\begin{problemblock}
\textbf{例 45.} 已知 \(n\) 阶矩阵 \(A,B,C\) 满足
\[
ABC=O,\qquad E\text{ 为单位矩阵}.
\]
记矩阵
\[
M_1=\begin{pmatrix}O&A\\BC&E\end{pmatrix},\qquad
M_2=\begin{pmatrix}AB&C\\O&E\end{pmatrix},\qquad
M_3=\begin{pmatrix}E&AB\\AB&O\end{pmatrix}
\]
的秩分别为 \(r_1,r_2,r_3\)，则（\quad）：
\[
\text{A. }r_1\le r_2\le r_3,\quad
\text{B. }r_1\le r_3\le r_2,\quad
\text{C. }r_3\le r_1\le r_2,\quad
\text{D. }r_2\le r_1\le r_3.
\]
\end{problemblock}
\begin{solutionblock}
\analysis{先求三个秩。对 \(M_1\)，右下角 \(E\) 可逆，用 Schur 补：
\[
r(M_1)=n+r(O-AE^{-1}BC)=n+r(-ABC)=n.
\]
故
\[
r_1=n.
\]
对 \(M_2\)，它是分块上三角矩阵，右下角 \(E\) 可逆，消去右上角后得
\[
r(M_2)=n+r(AB),
\]
故
\[
r_2=n+r(AB).
\]
对 \(M_3\)，左上角 \(E\) 可逆，用 Schur 补：
\[
r(M_3)=n+r(O-AB\cdot E^{-1}\cdot AB)=n+r\bigl((AB)^2\bigr).
\]
因此
\[
r_3=n+r\bigl((AB)^2\bigr).
\]
由于
\[
r\bigl((AB)^2\bigr)\le r(AB),
\]
并且 \(r\bigl((AB)^2\bigr)\ge0\)，所以
\[
n=r_1\le r_3\le r_2.
\]
故选 B。}
\answer{B；\(r_1\le r_3\le r_2\)}
\examnote{分块矩阵秩比较题，优先用可逆对角块的 Schur 补把秩写成 \(n+r(\text{补块})\)。}
\end{solutionblock}
